\documentclass[11pt,a4paper]{article}
\usepackage[margin=27mm]{geometry}
\usepackage{amsmath,amssymb,amsthm,booktabs,longtable,array}
\usepackage{fontspec,unicode-math}
\setmainfont{STIXGeneral}
\setmathfont{STIX Two Math}
\usepackage{hyperref}
\setlength{\emergencystretch}{2em}
\newtheorem{theorem}{Theorem}
\newtheorem{lemma}[theorem]{Lemma}
\newtheorem{corollary}[theorem]{Corollary}
\newtheorem{proposition}[theorem]{Proposition}
\newcommand{\PSL}{\operatorname{PSL}}
\newcommand{\PGL}{\operatorname{PGL}}
\newcommand{\tr}{\operatorname{tr}}
\newcommand{\Aut}{\operatorname{Aut}}
\newcommand{\ip}{\mathrm{IP}}
\title{Common theory for projective string markings}
\author{human\_compression project}
\date{11 October 2026}
\begin{document}
\maketitle
\begin{abstract}
We give one reduction for projective string markings, without assuming the
intersection property. Proper ends are controlled by the common centralizer
of two involutions. A separate finite theorem shows that, when both proper end groups are
$S_4$ or $A_5$ and both outer labels are odd, the marking gives precisely
the 11-cell, the 57-cell, or the 4-simplex. The local enumeration,
universal presentations and geometric conclusions are proved once and serve
both polytope classifications and the later maniplex application. A half-turn
extension theorem supplies the common construction for full and maximal ends.
Dickson's subgroup classification, standard projective group structure and
the group--polytope correspondence are the classical results used below.
\end{abstract}

\section{Markings and classical results}
A \emph{string marking} of a group $G$ is a generating ordered tuple
$(r_0,\ldots,r_{n-1})$ of distinct nonidentity involutions with
$r_ir_j=r_jr_i$ whenever $|i-j|>1$. Put
$G_I=\langle r_i:i\in I\rangle$, with $G_\varnothing=1$.
Its labels are $|r_0r_1|,\ldots,|r_{n-2}r_{n-1}|$.
The marking is a string C-marking if, in addition,
\begin{equation}\label{eq:ip}
 G_I\cap G_J=G_{I\cap J}\qquad(I,J\subseteq\{0,\ldots,n-1\}).
\end{equation}
For rank four write $H=\langle r_0,r_1,r_2\rangle$ and
$J=\langle r_1,r_2,r_3\rangle$ for its two ends. Throughout,
$L(q)=\PSL_2(q)$ and $P(q)=\PGL_2(q)$, for a prime power $q$.

We use the following classical results, in exactly the forms stated below.
Dickson's classification says that a proper subgroup of $L(q)$ or $P(q)$,
for odd $q$, is cyclic, dihedral, elementary abelian of characteristic order,
a subgroup of an affine group $V\rtimes k^\times$, one of $A_4,S_4,A_5$,
or a conjugate of a standard subfield group $L(k)$ or $P(k)$.
Only groups which actually embed are intended. In the subfield alternatives
$k\subseteq\mathbb F_q$ and the embedding is the natural matrix embedding,
not an arbitrary abstract group isomorphism. The possible same-field proper
subgroup $L(q)<P(q)$ must be included. We use
\[
 |L(q)|=\frac{q(q^2-1)}{(2,q-1)},\qquad |P(q)|=q(q^2-1),
\]
the simplicity of $L(q)$ for $q\ge4$, and $C_{P(q)}(L(q))=1$.
The small identifications are $L(2)=P(2)\cong S_3$, $L(3)\cong A_4$,
$P(3)\cong S_4$, $L(5)\cong A_5$ and $P(5)\cong S_5$.
These subgroup and structural facts are classical; subgroup multiplicity
formulas from the sources are not inputs.

We also use the correspondence between ordered string C-markings and
regular abstract polytopes, including the following clauses. Rank-$i$ faces
are cosets of $G_{\{0,\ldots,n-1\}\setminus\{i\}}$; reversal gives duality;
a marking-preserving epimorphism of string C-groups induces a covering.
The local markings $[3,5]_5$ and $[5,3]_5$ below are the usual antipodal
icosahedral and dodecahedral quotients. These names of the local maps, and the identification of the simplex,
are standard geometric interpretations.
They do not replace the required verification of the actual markings.

All groups $L(q),P(q)$ are centre-free and directly indecomposable. For the
simple groups this is immediate. To see the assertion for $P(q)$ when
$q\ge4$, let commuting subgroups $A,B$ generate $P(q)$; both are normal.
The subgroup $[L(q),A]$ lies in $L(q)\cap A$ and is normal in the simple
group $L(q)$. If it is nontrivial, $L(q)\le A$, so $B$ centralizes $L(q)$
and is trivial. Otherwise $A$ centralizes $L(q)$ and is trivial. The small
cases follow from the normal-subgroup lists
$1,A_3,S_3$; $1,V_4,A_4$; and $1,V_4,A_4,S_4$.
In any string marking of these groups every adjacent product has order
at least three. Indeed, a commuting adjacent pair splits the generators
into two nontrivial commuting generated subgroups. Their intersection is
central, hence trivial, and the resulting direct product contradicts
indecomposability. This argument does not assume \eqref{eq:ip}.

\section{Common centralizers}
In a string marking, the fourth generator commutes with the first two.
Their common centralizer therefore controls both the possible rank and
the choice of a fourth generator.
\begin{lemma}\label{lemma:centralizer}
In characteristic two, the centralizer of a nonidentity involution in
$P(q)$ is abelian. Let $k\subseteq K$ be finite fields of odd characteristic,
and let $a,b\in P(k)$ be involutions with $|ab|>2$. Their common centralizer
in $P(K)$ lies in $P(k)$. It is trivial when $ab$ is unipotent and has
order two otherwise. In the latter case its nonidentity element is represented
by the invertible matrix
\begin{equation}\label{eq:trace}
 Z=A-\tfrac12\tr(A)I,
\end{equation}
where $A\in\operatorname{GL}_2(k)$ represents $ab$.
\end{lemma}
\begin{proof}
In characteristic two an involution is unipotent. Over an algebraic closure
conjugate it to $x\mapsto x+1$. Its unique fixed point is infinity, so a
commuting projectivity is $x\mapsto cx+d$. Commutation gives $c=1$;
these translations form an abelian group. The centralizer over the original
field is a subgroup of this group.

In odd characteristic first suppose $u=ab$ is semisimple. Over the algebraic
closure write $u:x\mapsto\lambda x$, with $\lambda\ne\pm1$.
A commuting projectivity preserves each of its two fixed points, since
interchanging them would send $u$ to $u^{-1}\ne u$. It is therefore
$x\mapsto cx$. As $aua^{-1}=u^{-1}$, the involution $a$ has the form
$x\mapsto d/x$. Commutation with $a$ gives $c^2=1$, so the common
centralizer is $\{1,x\mapsto-x\}$. If $A$ has eigenvalues
$\alpha,\beta$, the matrix in \eqref{eq:trace} has the nonzero eigenvalues
$(\alpha-\beta)/2$ and $(\beta-\alpha)/2$. It represents $x\mapsto-x$
and has entries in $k$. Rescaling $A$ only rescales $Z$, proving the
claimed descent, including the nonsplit case.

If $u$ is unipotent, conjugate it to $x\mapsto x+1$. Its centralizer again
consists of translations $x\mapsto x+c$. An involution conjugating it to
its inverse has the form $x\mapsto-x+d$. A commuting translation satisfies
$2c=0$, and hence $c=0$. These cases exhaust $u$.
\end{proof}

\begin{lemma}\label{lemma:rank}
Let $G$ be centre-free with a string marking. Its first two generators do
not commute. If every noncommuting pair of involutions has at most one
nonidentity common centralizer element, the rank is at most four. If every
involution centralizer is abelian, the rank is at most three.
\end{lemma}
\begin{proof}
If $r_0$ commuted with $r_1$, it would commute with every generator by the
string relations, so would be a nonidentity central element. At rank at
least five, both $r_3$ and $r_4$ centralize $r_0,r_1$, contradicting the
first centralizer hypothesis and their distinctness. Under the second
hypothesis a rank-four marking puts $r_0,r_1$ in the abelian centralizer
of $r_3$, again a contradiction.
\end{proof}

\begin{lemma}\label{lemma:odd}
Let a rank-four string marking lie in $P(q)$, with $q$ odd and adjacent
labels greater than two. If $H$ is proper, then $|r_0r_1|$ is odd; if
$J$ is proper, then $|r_2r_3|$ is odd.
\end{lemma}
\begin{proof}
The nonidentity element $r_3$ centralizes $D=\langle r_0,r_1\rangle$.
Lemma~\ref{lemma:centralizer} excludes the unipotent case and gives a
unique nonidentity common centralizer element. If $m=|r_0r_1|$ were even,
$(r_0r_1)^{m/2}$ would be that element. Then $r_3\in H$, contrary to
$G=\langle H,r_3\rangle$ and properness. Reverse the indices for $J$.
\end{proof}

\begin{theorem}[Projective reduction]\label{thm:projective}
A string marking of $L(q)$ or $P(q)$ has rank at most four; rank four
requires odd $q$. In rank four every proper end of $L(q)$ is $S_4$ or
$A_5$. Every proper end of $P(q)$ is $S_4,A_5$, or its natural subgroup
$L(q)$. If both ends are proper, both are $S_4$ or $A_5$ and both outer
labels are odd. These statements require no intersection property.
\end{theorem}
\begin{proof}
The rank claims follow from Lemmas~\ref{lemma:centralizer} and
\ref{lemma:rank}, since the projective groups are centre-free. This also
handles $q=2$; $L(3)$ has no generating tuple of involutions, since its
involutions generate only $V_4$. All adjacent labels in the remaining
cases exceed two by the first section.

Let $H=\langle r_0,r_1,r_2\rangle$ be proper and put
$D=\langle r_0,r_1\rangle$. A cyclic or elementary abelian group cannot
contain the nonabelian group $D$. Involutions in $A_4$ generate only
$V_4$. In an affine group in odd characteristic every involution is
$x\mapsto-x+v$; two such involutions commute only when they coincide.
Thus the distinct commuting endpoints $r_0,r_2$ exclude this alternative,
including affine dihedral groups of order $2p$.

If $H$ is dihedral, its distinct commuting involutions $r_0,r_2$ force
its rotation subgroup to have even order. (In an odd dihedral group two
reflections commute only when equal.) Its central involution $z$ and
$r_3$ both centralize $D$. Lemma~\ref{lemma:centralizer} gives $r_3=z$,
so $r_3\in H$, a contradiction. This exclusion uses properness, not an
intersection identity inside $H$.

If $H$ is a natural subfield group, conjugate it into $P(k)$.
Lemma~\ref{lemma:centralizer} puts $r_3$ in the same $P(k)$, hence puts
all of $G$ there. For a proper subfield of size $s$ we have $q\ge s^2$,
and
\[
 |G|\ge q(q^2-1)/2>s(s^2-1)=|P(k)|,
\]
a contradiction. The only remaining same-field proper alternative is
$H=L(q)<P(q)$. This proves the asserted lists of proper ends. The
qualification ``natural'' is essential: exceptional $A_5\cong L(5)$
and $S_4\cong P(3)$ in other characteristics are not subfield exclusions.

Suppose now that both ends are proper. Lemma~\ref{lemma:odd} makes both
outer labels odd. If $H=L(q)<P(q)$, then $H$ is normal of index two;
$r_2\in H$ and $r_3\notin H$. The image of $r_2r_3$ in $P(q)/H$
therefore has order two, impossible for its odd order. Reversal excludes
the same alternative for $J$. The remaining ends are $S_4,A_5$.
\end{proof}

\section{End groups isomorphic to \texorpdfstring{$S_4$ or $A_5$}{S4 or A5}}
For a triple $f=(m,n,k)$ write
\begin{equation}\label{eq:local}
 [m,n]_k=\langle x,y,z\mid x^2=y^2=z^2=(xz)^2=
 (xy)^m=(yz)^n=(xyz)^k=1\rangle.
\end{equation}
The third parameter is the Petrie order, not an adjacent label.
The six ordered local classes are
\begin{equation}\label{eq:localtypes}
\begin{array}{c|ccc}
 S_4&(3,3,4)&(3,4,3)&(4,3,3)\\
 A_5&(3,5,5)&(5,3,5)&(5,5,3).
\end{array}
\end{equation}
For $f=(m,n,k)$ and $v=(n,s,l)$ in this list define
\begin{equation}\label{eq:universal}
\begin{split}
 U(f,v)=\langle r_0,r_1,r_2,r_3\mid{}
 &r_i^2=1,\ (r_ir_j)^2=1\ (j-i>1),\\
 &(r_0r_1)^m=(r_1r_2)^n=(r_2r_3)^s=1,\\
 &(r_0r_1r_2)^k=(r_1r_2r_3)^l=1\rangle.
\end{split}
\end{equation}
Matching the middle label gives fourteen ordered pairs. Reversal sends
$(f,v)$ to $(v^*,f^*)$, where $(a,b,c)^*=(b,a,c)$.
There are ten reversal orbits, six fixed. Their complete table is:
\begin{center}\small
\begin{tabular}{rcccccl}
\toprule
Row&$f$&$v$&Order&Group&C-marking&Outer labels\\
\midrule
1&$(5,5,3)$&$(5,5,3)$&1&trivial&no&$5,5$\\
2&$(5,5,3)$&$(5,3,5)$&1&trivial&no&$5,3$\\
3&$(5,3,5)$&$(3,5,5)$&3420&$L(19)$&yes&$5,5$\\
4&$(5,3,5)$&$(3,4,3)$&60&$A_5$&no&$5,4$\\
5&$(5,3,5)$&$(3,3,4)$&1&trivial&no&$5,3$\\
6&$(4,3,3)$&$(3,4,3)$&96&$2^4:S_3$&yes&$4,4$\\
7&$(4,3,3)$&$(3,3,4)$&24&$S_4$&no&$4,3$\\
8&$(3,5,5)$&$(5,3,5)$&660&$L(11)$&yes&$3,3$\\
9&$(3,4,3)$&$(4,3,3)$&1&trivial&no&$3,3$\\
10&$(3,3,4)$&$(3,3,4)$&120&$S_5$&yes&$3,3$\\
\bottomrule
\end{tabular}
\end{center}
Here an entry ``yes'' means that the specified generators, not merely some
marking of the abstract group, form a string C-group. Rows which collapse
do not have the indicated local groups; the local triples label their
presentation input. The order-$96$ row has adjacent outer labels $4,4$,
not $3,3$.

\begin{lemma}[Exact finite certificate]\label{lemma:finite}
Every rank-three string marking of $S_4$ or $A_5$ has the intersection
property. Up to the full automorphism group its ordered type and Petrie
order are exactly \eqref{eq:localtypes}; \eqref{eq:local} is an exact
presentation in each case. The ten-row table above is correct, including
its groups, specified markings, collapse assertions and reversal coverage.
In rows $3,8,10$, reversal extends to an automorphism of the marked group.
\end{lemma}
\begin{proof}
The complete program \texttt{required\_finite.g} enumerates all involutions
in $S_4,A_5$, all ordered pairs of distinct commuting endpoints $x,z$,
and every middle involution $y$. Its admission tests are only distinctness
and generation. It admits $72$ and $360$ triples respectively. For every
admitted triple it then checks both adjacent orders and all $64$
subset-pair identities in \eqref{eq:ip}. The intersection property is therefore proved for every admitted string
marking; it is not assumed during the enumeration. The full
automorphism orbits give the three classes in each row of
\eqref{eq:localtypes}. The code constructs each literal presentation
\eqref{eq:local} and computes its order $24$ or $60$. The marking
provides a surjection onto the corresponding finite group; equality of
orders makes it an isomorphism.

The code independently generates all fourteen compatible pairs, partitions
by reversal and checks the ten supplied representatives and six fixed
pairs. For each literal presentation \eqref{eq:universal}, it computes
the finite presentation order, constructs an isomorphic permutation group
and transports the four original generators. It checks generator orders and distinctness, then tests IP over all sixteen
subsets. Every accepted marking satisfies all $256$ subset-pair identities;
a rejected row may fail a generator condition or terminate at its first
failed intersection. Thus the negative entries certify failure of the
complete string C-marking conditions, not necessarily a separate IP test.
Rows $3,8,10$ are identified by actual group isomorphisms with
$L(19),L(11),S_5$; generator reversal is checked by constructing the
homomorphism and testing bijectivity. In rows $4,7$ it verifies the indicated group and that
one end is the whole group. In row $6$ it verifies an elementary abelian
normal subgroup of order $16$, quotient $S_3$ and a complement. The four
trivial rows are checked directly. These exact computations prove the
finite assertions using GAP's algorithms for finite groups and finitely
presented groups; they are not Lean kernel computations.
\end{proof}

\begin{theorem}[Exceptional ends and odd outer labels]\label{thm:finite}
Let $(G;r_0,r_1,r_2,r_3)$ be a string marking whose two ends are proper
and isomorphic to $S_4$ or $A_5$. If both outer labels are odd, then,
up to isomorphism preserving the ordered marking, it is exactly one of
rows $3,8,10$. Each row supplies one class, satisfies IP, and has type
$\{5,3,5\}$, $\{3,5,3\}$ or $\{3,3,3\}$ respectively.
In particular the statement does not require $G$ to be a projective group.
\end{theorem}
\begin{proof}
The local census and its exact presentations give a marking-preserving
surjection from one of the ten groups $U(f,v)$ to $G$, after possibly
reversing the marking. It is necessary to examine quotients, not just
the universal groups themselves. The trivial rows cannot give nonidentity
markers. In rows $4,7$, an end already equals the full universal group;
this equality survives every quotient, contradicting properness. Row $6$
has first and last local product orders four, whereas the given ends
have their exact local markings and hence these exact orders. It is
excluded by oddness.

The remaining universals are the simple groups $L(19),L(11)$ and $S_5$.
The simple groups have no nontrivial proper quotient. The normal subgroups
of $S_5$ are $1,A_5,S_5$, so its only nontrivial proper quotient is $C_2$,
which cannot have an end $S_4$ or $A_5$. Thus every surviving map is an
isomorphism on all four distinguished generators. Lemma~\ref{lemma:finite}
gives IP and the asserted types. The local classes were taken under full
automorphisms, so no marking was lost through a conjugacy restriction.
The reversal automorphisms verified below identify the two representatives
of each reversal orbit. Consequently each surviving row gives one class under isomorphisms
preserving the ordered generators.
\end{proof}

\noindent\emph{Sharpness of the parity hypothesis.} The verified order-$96$
row has proper $S_4$ ends, type $\{4,3,4\}$ and IP. It therefore
contradicts the conclusion of three cases if the requirement that both
outer labels be odd is removed. We make no classification of all quotients without that
hypothesis.

\begin{corollary}\label{cor:projective}
Every string marking of $L(q)$ or $P(q)$ of rank at least four with both
ends proper is one of these three markings. The first two occur in
$L(19),L(11)$ respectively; the third occurs in $P(5)$.
A string C-marking automatically has both ends proper, so the same
conclusion applies to all projective string C-markings of rank at least four.
\end{corollary}
\begin{proof}
Theorem~\ref{thm:projective} supplies rank four, the exceptional ends and
odd outer labels; apply Theorem~\ref{thm:finite}. The order formulas are
strictly increasing in odd $q$. The groups $L(11),L(19)$ are simple,
whereas odd-characteristic $P(q)$ has a proper subgroup of index two.
The group $S_5$ is not simple and has order $120=|P(5)|$; the small
$L(3)$ is not $S_5$. These observations give exactly the stated field
values and projective types. Finally, if an omitted generator belonged
to its complementary subgroup, \eqref{eq:ip} would intersect that subgroup
with its generator subgroup nontrivially. Hence C-markings have proper ends.
\end{proof}

\section{Geometry and universal coverings}
\begin{theorem}\label{thm:geometry}
The three markings in Theorem~\ref{thm:finite} give the 57-cell, the
11-cell and the 4-simplex. Their types, local maps and face vectors are
\[
\begin{array}{c|c|c|c}
\text{group}&\text{type}&\text{facet / vertex figure}&(f_0,f_1,f_2,f_3)\\
L(19)&\{5,3,5\}&\text{hemi-dodecahedron / hemi-icosahedron}&(57,171,171,57)\\
L(11)&\{3,5,3\}&\text{hemi-icosahedron / hemi-dodecahedron}&(11,55,55,11)\\
S_5&\{3,3,3\}&\text{tetrahedron / tetrahedron}&(5,10,10,5).
\end{array}
\]
All three are self-dual and universal among regular polytopes with their
respective facet and vertex-figure types. The 11-cell edge graph is
$K_{11}$. The even-word subgroup is the full group in the two simple
rows, and is $A_5$ in the simplex row.
\end{theorem}
\begin{proof}
The two simple rows have local presentations $[3,5]_5,[5,3]_5$, with the
standard antipodal map interpretations. Their exact global presentations
are the string Coxeter relations of types $\{3,5,3\}$ and $\{5,3,5\}$,
together with
\[
 (r_0r_1r_2)^5=(r_1r_2r_3)^5=1.
\]
For the third row send the four generators to $(12),(23),(34),(45)$.
These generate $S_5$ and satisfy every relator in
\eqref{eq:universal}, including the two fourth powers of consecutive
triple products. Both groups have order $120$, so this is an isomorphism
of marked groups with the standard simplex marking.

The program verifies that generator reversal extends to an automorphism
in each of the three rows; the group--polytope correspondence gives self-duality.
It also computes each of the four complementary-subgroup indices, giving
the displayed face vectors. In the simple rows this has the following
direct interpretation. Both ends have order $60$. An edge stabilizer is
$\langle r_0\rangle\times\langle r_2,r_3\rangle$, since the factors
commute and IP makes their intersection trivial. Its order is $4m$ for
outer label $m=3$ or $5$; the two-face stabilizer has the same order.
Thus $660/60=11$, $660/12=55$, $3420/60=57$ and $3420/20=171$ give the
first two vectors. For the 11-cell the code constructs actual vertex and
edge cosets and checks their incidences: every edge has two endpoints and
the $55$ resulting unordered pairs are distinct. There are exactly $55$
pairs on eleven vertices, proving $K_{11}$ rather than inferring it merely
from face counts.

The even-word subgroup is normal of index at most two: in its quotient
all involutions have the same image. Simplicity makes it the full group
in the two simple rows. In the simplex marking parity is the ordinary
sign homomorphism, giving $A_5$.

Let $Q$ be any regular polytope with one of the three indicated pairs of
local maps. Its distinguished generators satisfy the corresponding
literal presentation \eqref{eq:universal}. There is therefore a marked
epimorphism from the already verified string C-group of that row to its
automorphism group. The classical covering theorem gives a covering of
$Q$ by the corresponding polytope above. No finiteness assumption on $Q$
is needed: the epimorphism itself forces its group to be finite. This is
universality of these exact markings, not an assumption that every
abstract amalgam has IP.
\end{proof}

\section{The common extension construction}
We use the standard torus structure in $L(q)$ for odd $q$: its split and
nonsplit cyclic tori have orders $(q-1)/2,(q+1)/2$, with dihedral
normalizers of twice these orders. The split torus consists of diagonal
projectivities of square eigenvalue ratio. The nonsplit torus comes from
multiplication on $\mathbb F_{q^2}$ modulo scalars, with the determinant-
square condition. Reflections in the normalizer interchange the two
eigenlines. We use these classical torus facts, without assuming maximality of
these normalizers.

\begin{theorem}[Half-turn extension]\label{thm:extension}
Let $q$ be odd, choose $\alpha\in\{1,-1\}$ with $q\equiv\alpha\pmod4$,
and put $N=(q-\alpha)/2$. Let $(a,b,c)$ be a rank-three string marking
of $H\le L(q)$, with $|bc|>2$ and $|ab|=m>1$ an odd divisor of $N$.
There is a unique nonidentity projectivity $z$ centralizing $a,b$;
it lies in $L(q)$ and $(a,b,c,z)$ is a string marking of $\langle H,z\rangle$.
If $H=L(q)$, this group remains $L(q)$. If $H$ is maximal proper in
$L(q)$ and $C_H(\langle a,b\rangle)=1$, it is again $L(q)$ and its
first end is exactly $H$.
\end{theorem}
\begin{proof}
The element $ab$ is semisimple, since $m$ divides $q-\alpha$ and hence
is prime to the characteristic. It has two eigenlines over at most a
quadratic extension. The two torus orders are coprime consecutive integers;
since $m>1$ divides $N$, its torus has order $N$. Both $a,b$ invert
$ab$, so preserve its unordered eigenline pair and normalize this torus.
The even cyclic torus has a unique involution $z$. It commutes with the
torus and its inverting normalizer coset, so centralizes $a,b$.
Lemma~\ref{lemma:centralizer} proves its asserted uniqueness even in
$P(q)$.

The dihedral group $\langle a,b\rangle$ has odd rotation order and
trivial centre. Thus $z$ does not belong to it and differs from $a,b$.
It also differs from $c$, because it commutes with $b$ whereas $|bc|>2$.
The four generators are distinct nonidentity involutions and have every
required distant commutation. They generate $\langle H,z\rangle$ by
construction. If $H=L(q)$ there is nothing further to prove. Under the
second hypothesis $z\notin H$, because its membership would give a
nonidentity element of the stated trivial internal centralizer.
Maximality then gives $\langle H,z\rangle=L(q)$.
\end{proof}
The extension theorem concerns string markings; it does not assert IP
for the extended object. The maniplex application proves the correspondence with regular
maniplexes and hence shows that the first end describes the original map.
It proves its constructions for every prime in the stated ranges and
identifies the generators with those in the source by explicit conjugacy.
These conclusions do not rely on a sample of finite cases.

\section{Attribution and verification}
The six local classes, ten universal groups and the projective polytope
classifications are prior results of Leemans--Schulte \cite{LS1,LS2}; the
11-cell and 57-cell are the prior constructions of Gr\"unbaum and Coxeter.
The rank bound and classification of proper end groups, proved without
assuming the intersection property, and the torus extension recover
results of Leemans--Toledo \cite{LT}.
Our common proof combines projective centralizer descent with a finite
theorem for proper end groups $S_4$ and $A_5$ with odd outer labels.
That theorem derives the exact marked groups, their intersection property
and uniqueness from one calculation. It supplies the common argument for
the polytope and maniplex applications. We make no claim to be the first
to use these proof methods. The classical subgroup types, embeddings,
torus structure and geometric correspondences remain attributed to the
previous literature.

The finite calculations rely on the complete program
\texttt{required\_finite.g}. Besides the local and universal calculations
proved above, it records the exact $A_5$ involution-centralizer, $L(23)$
symmetric-end and self-dual $L(13)$ witnesses used for the literal source
corrections in the third application. Their exact conclusions are stated
there. The optional \texttt{validation.g} contains finite samples of
uniform constructions and controls against weakening hypotheses. Those
samples are not premises of the all-prime proofs, which are written in
full in the application. No mathematical branch of the required program
is excluded from its cost.

Existing Lean components verify the normalizer parity argument, the
abstract string centralizer rank obstruction, matrix normal-form identities
and a generic extension interface. They do not formalize Dickson's theorem,
projective descent in full, the GAP table, torus structure or the geometric
correspondence. The source ledger records this boundary. A successful
compilation of those components is not a Lean proof of this whole paper.
The companion applications state exactly which previous results they
recover and which additional statements they leave aside. This edition changes proof sharing, not the
number of sources or questions credited.

\begin{thebibliography}{9}
\bibitem{LS1} D. Leemans and E. Schulte, \emph{Groups of type $L_2(q)$ acting
on polytopes}, arXiv:math/0606660v1, bound fourteen-page author edition;
Advances in Geometry 7 (2007), 529--539.
\bibitem{LS2} D. Leemans and E. Schulte, \emph{Polytopes with groups of type
$\PGL_2(q)$}, arXiv:0909.1991v1, bound nine-page edition.
\bibitem{LT} D. Leemans and M. Toledo, \emph{Maniplexes with automorphism
group $\PSL(2,q)$}, bound eleven-page author manuscript dated 15 May 2023.
\bibitem{MS} P. McMullen and E. Schulte, \emph{Abstract Regular Polytopes},
Cambridge University Press, 2002; string C-group and covering correspondence.
\end{thebibliography}
\end{document}
